\begin{solapartat}
\begin{minipage}[t]{0.5\linewidth}
\punts{0,3} $E = hf - \underbrace{hf_0}_{W_\textit{extracció}} = h(f - f_0)$

\bigskip
A partir del pendent de la recta es calcula la constant de Planck.

\bigskip
Triem per exemple, els punts:\\
$(1,5\cdot 10^{15}, 2,07)$ i $(1,0\cdot 10^{15}, 0)$
\end{minipage}\hfill
\begin{minipage}[t]{0.46\linewidth}
\vspace{0pt}
\figura[1]{sol_fig1.png}
\end{minipage}

\punts{0,7}
\[ \mathit{pendent} \approx \frac{(2,07\times 1,6\times 10^{-19}) - 0}{1,5\times 10^{15} - 10^{15}} = 6,62\times 10^{-34}\ \mathrm{Js} \]

Es poden acceptar valors aproximats ja que estem llegint una gràfica.
\end{solapartat}

\begin{solapartat}
\punts{0,4} La freqüència llindar $f_0$, és, segons la gràfica: $10^{15}$~Hz

La energia mínima d'extracció dels electrons és:

\punts{0,6} $W_\textit{extracció} = hf_0 = 6,62\times 10^{-34}\times 10^{15} = 6,62\times 10^{-19}\ \mathrm{J} = 4,14\ \mathrm{eV}$
\end{solapartat}
